\section*{Bab \ref{ch:b2:pt-paths} --- Termobarometri dan Lintasan P--T--t}

\begin{solution}{exo:b2:pt-paths:1}
(a) Sampel A terletak di zona klorit, B di zona biotit, C di zona garnet, D di zona staurolit, E di zona kianit,
dan F di zona silimanit. (b) Lima: isograd biotit, garnet, staurolit, kianit, dan silimanit,
masing-masing di antara dua sampel yang berurutan, dengan derajat naik dari A ke F. (c) Isograd
hanya mengatakan bahwa suatu reaksi telah berlangsung; gradien lapangan memerlukan tekanan dan suhu, yang
berasal dari reaksi yang terkalibrasi atau dari termobarometri \omterm{def:b2:metamorphic-facies:assemblage}{himpunan mineral} puncak setiap
batuan.
\end{solution}

\begin{solution}{exo:b2:pt-paths:2}
(a) $(X_{\text{Mg}}/X_{\text{Fe}})_{\text{grt}} = 0.10/0.80 = 0.125$,
$(X_{\text{Mg}}/X_{\text{Fe}})_{\text{bt}} = 0.45/0.50 = 0.900$; $K_D = 0.1389$, $\ln K_D =
-1.9741$. (b) $3R = 24.942$; penyebut $20.19 + 24.942\times1.9741 = 69.428$; pada
\qty{5}{kbar}, $T = 54185/69.428 = \qty{780.5}{K} = \qty{507}{\celsius}$; pada
\qty{7}{kbar}, $T = (52630 + 2177)/69.428 = \qty{789.4}{K} = \qty{516}{\celsius}$.
(c) Dua kilobar hanya mengubahnya sebesar \qty{9}{K}: sebuah termometer.
\end{solution}

\begin{solution}{exo:b2:pt-paths:3}
(a) $z = P/(\bar\rho g) = 6.0\times10^8/(2800\times9.81) = \qty{21.8}{km}$.
(b) $(580 - 10)/21.8 = \qty{26}{\celsius\per\km}$. (c) Pada \qty{580}{\celsius} garis itu
terletak pada $3.67 + 0.8\times2.10 = \qty{5.35}{kbar}$; pada \qty{6.0}{kbar} batuan berada di
sisi tekanan tinggi: kianit.
\end{solution}

\begin{solution}{exo:b2:pt-paths:4}
(a) Pereaksi $2\times140.693 + 3\times18.015 = \qty{335.431}{g}$; produk $277.112 +
58.320 = \qty{335.432}{g}$: setara sampai batas pembulatan. (b) $2\times43.65 =
\qty{87.30}{cm^3}$ menjadi $107.50 + 24.63 = \qty{132.13}{cm^3}$: $+\qty{51.4}{\percent}$.
(c) Olivin $281.386/87.30 = \qty{3.223}{g.cm^{-3}}$; produk $335.432/132.13 =
\qty{2.539}{g.cm^{-3}}$. (d) Batuan itu mengembang sebesar setengah volumenya dan menjadi lebih ringan; batuan itu
juga menyimpan air di dalam mineral-mineralnya.
\end{solution}

\begin{solution}{exo:b2:pt-paths:5}
(a) Rata-ratanya $-1.206$ pada \qty{1072.15}{K} dan $-1.782$ pada \qty{823.15}{K}. Dengan $1/T =
9.3271\times10^{-4}$ dan $1.21485\times10^{-3}$, $a = (-1.206 + 1.782)/(9.3271\times10^{-4}
- 1.21485\times10^{-3}) = -2041.5$ K dan $b = -1.206 + 2041.5/1072.15 = 0.6982$.
(b) $\Delta_r H + 2070\times0.311 = 2041.5\times24.942 = \qty{50920}{J}$, sehingga $\Delta_r H =
\qty{50276}{J}$; $\Delta_r S = 0.6982\times24.942 = \qty{17.41}{J.K^{-1}}$. (c) Pada
\qty{5}{kbar}: untuk $K_D = 0.18$, 588.1 terhadap \qty{587.5}{\celsius}; untuk $0.10$, 419
terhadap \qty{425}{\celsius}; untuk $0.30$, 819 terhadap \qty{806}{\celsius}. Di dalam rentang
percobaan, kedua kalibrasi itu sesuai, walaupun konstanta-konstantanya berbeda sebesar
\qty{2.4}{kJ} dan \qty{2.8}{J.K^{-1}}; galat $\Delta_r H$ dan $\Delta_r S$
saling mengimbangi di sana dan membesar di luarnya. Suatu kalibrasi tidak boleh dipakai jauh
di luar kondisi percobaannya.
\end{solution}

\begin{solution}{exo:b2:pt-paths:6}
(a) Pertukaran: $(113.12 + 154.30) - (115.32 + 149.65) = \qty{+2.45}{cm^3} =
\qty{0.245}{J.bar^{-1}}$, terhadap 0.311. GASP: $125.28 + 2\times44.15 + 22.69 -
3\times100.79 = \qty{-66.10}{cm^3} = \qty{-6.61}{J.bar^{-1}}$, terhadap $-6.21$, nilai
kalibrasi pada \qty{800}{\celsius} dan \qty{10}{kbar}, tempat pemuaian termal dan
pemampatan mengubah volumenya. Dalam kedua kasus, perubahan GASP dua puluh sampai tiga puluh kali
perubahan pertukaran. (b) $\partial T/\partial P = 0.311/62.96 = \qty{4.9e-3}{K/bar}$:
\qty{4.9}{K} per kilobar. (c) Kemiringan garisnya $22.80 + 8.314\times3\ln 0.20/6.21
= 22.80 - 6.46 = \qty{16.3}{bar/K}$; \qty{17}{K} menggeser tekanan sebesar
\qty{280}{bar}.
\end{solution}

\begin{solution}{exo:b2:pt-paths:7}
(a) $P_0 = 22.80\times600 - 1093 = \qty{12587}{bar}$; $3\ln(0.08/0.30) = -3.9653$; $RT =
8.314\times873.15 = \qty{7259.4}{J}$; $P = 12587 + 7259.4\times(-3.9653)/6.21 = 12587 -
4635 = \qty{7952}{bar}$. (b) $3\ln(0.08/0.35) = -4.4277$, $P = 12587 - 5176 =
\qty{7411}{bar}$: \qty{0.54}{kbar} lebih rendah, jadi plagioklas harus dianalisis dengan cermat.
(c) $7952\times10^5/(2800\times9.81) = \qty{29.0}{km}$.
\end{solution}

\begin{solution}{exo:b2:pt-paths:8}
(a) $z = P/(\bar\rho g)$: \qty{118.9}{km}, \qty{34.0}{km}, dan \qty{15.3}{km}.
(b) $(118.9 - 34.0)/2.2 = \qty{38.6}{km/Ma} = \qty{3.9}{cm.yr^{-1}}$ dan $(34.0 -
15.3)/1.1 = \qty{17.0}{km/Ma} = \qty{1.7}{cm.yr^{-1}}$, sedikit di atas nilai terbitan
3.4 dan 1.6: konversi tekanan ke kedalaman yang mereka pakai menghasilkan selisih kedalaman yang sedikit
lebih kecil. (c) $\sigma = \sqrt{0.9^2 + 0.9^2} = \qty{1.3}{Ma}$, jadi $2.2 \pm 1.3$
juta tahun dan laju antara $84.9/3.47 = 24$ dan $84.9/0.93 = \qty{91}{km/Ma}$,
2.4 sampai \qty{9.1}{cm.yr^{-1}}: bagaimanapun juga cepat. (d) Laju terendah pun sekitar lima puluh kali
\qty{0.5}{mm.yr^{-1}} model itu: erosi saja tidak mungkin melakukannya.
\end{solution}

\begin{solution}{exo:b2:pt-paths:9}
(a) $\partial T/\partial\ln K_D = 3RT/(\Delta_r S - 3R\ln K_D) = 24.942\times860.6/62.96 =
\qty{341}{K}$; $\partial T/\partial P = 0.311/62.96 = \qty{4.9}{K}$ per kilobar.
(b) $\sigma_T = \sqrt{(341\times0.05)^2 + 4.9^2} = \sqrt{17.0^2 + 4.9^2} =
\qty{17.7}{K}$. (c) Galat analitis pada $K_D$ yang dominan; galat ini saja sudah dua pertiga
dari \qty{25}{\celsius} kalibrasi terbaik, yang juga mencakup galat
kalibrasi itu sendiri.
\end{solution}

\begin{solution}{exo:b2:pt-paths:10}
(a) Garnet $0.10/0.75 = 0.1333$, biotit $0.42/0.52 = 0.8077$, $K_D = 0.1650$, $\ln K_D =
-1.8018$, penyebut $20.19 + 24.942\times1.8018 = 65.130$; $3\ln(0.07/0.25) = -3.8189$.
Pada \qty{6000}{bar}: $T = 54496/65.130 = \qty{836.7}{K} = \qty{563.6}{\celsius}$, lalu $P
= 22.80\times563.6 - 1093 + 8.314\times836.7\times(-3.8189)/6.21 = \qty{7479}{bar}$.
Pada \qty{7479}{bar}: $T = \qty{570.6}{\celsius}$, $P = \qty{7603}{bar}$. Pada
\qty{7603}{bar}: $T = \qty{571.2}{\celsius}$, $P = \qty{7614}{bar}$; sekali lagi, $T =
\qty{571.3}{\celsius}$. Batuan itu mencapai kesetimbangan pada sekitar \qty{571}{\celsius} dan
\qty{7.6}{kbar}. (b) Pada \qty{571}{\celsius} garis kianit $=$ silimanit terletak pada $3.67
+ 0.71\times2.10 = \qty{5.16}{kbar}$: kianit. (c) Barometer itu dikalibrasi untuk
kianit; dengan silimanit, reaksinya, dan garisnya, berbeda, karena kedua polimorf itu tidak
memiliki energi Gibbs yang sama.
\end{solution}

\begin{solution}{exo:b2:pt-paths:11}
(a) $\tau = L^2/(\pi^2\kappa)$: untuk \qty{35}{km}, $1.225\times10^9/(9.870\times10^{-6}) =
\qty{1.24e14}{s} = \qty{3.9}{Ma}$; untuk \qty{70}{km}, empat kali lebih lama,
\qty{15.7}{Ma}. (b) Tujuh puluh juta tahun adalah empat sampai delapan belas kali waktu-waktu itu: panas
berpindah menembus kerak yang menebal jauh lebih cepat daripada erosi mengupasnya. (c) Batuan itu
mencapai kedalaman terbesarnya seketika, saat penumpukan; panas datang selama puluhan juta tahun
berikutnya, sementara erosi perlahan-lahan mengangkatnya: suhunya terus
naik untuk sementara waktu selama dekompresi. (d) Erosi dalam tujuh juta tahun akan
mengangkat batuan itu sebelum konduksi sempat memanaskannya: puncak termal yang lebih rendah, yang tercapai
lebih awal, pada tekanan yang lebih tinggi, dan lintasan dingin kembali ke permukaan.
\end{solution}

\begin{solution}{exo:b2:pt-paths:12}
(a) Kerak memberikan beban $35000\times2800\times9.81 = \qty{0.961}{GPa}$; sisa
\qty{1.839}{GPa} memerlukan $1.839\times10^9/(3300\times9.81) = \qty{56.8}{km}$ mantel: kedalaman
\qty{91.8}{km}, sesuai dengan \qty{90}{km} atau lebih pada studi aslinya. (b) $(4.0 -
0.961)\times10^9/(3300\times9.81) = \qty{93.9}{km}$, kedalaman \qty{128.9}{km}.
(c) $700/91.8 = \qty{7.6}{\celsius\per\km}$: rasio tekanan terhadap suhu yang sangat tinggi,
seri zona subduksi. (d) $(91.8 - 35)/34 = \qty{1.7}{Ma}$.
\end{solution}

\begin{solution}{pb:b2:pt-paths:1}
\noindent\textbf{Bagian I.}

\textbf{1.} Inti garnet $0.080/0.760 = 0.1053$, biotit $0.40/0.58 = 0.6897$; $K_D =
0.1527$, $\ln K_D = -1.8793$.

\textbf{2.} Tepi $0.120/0.720 = 0.1667$, biotit matriks $0.46/0.53 = 0.8679$; $K_D =
0.1921$, $\ln K_D = -1.6497$.

\textbf{3.} $3R = 24.942$; pembilang $52630 + 0.311\times6000 = 54496$; penyebut $20.19
+ 24.942\times1.8793 = 67.064$; $T = \qty{812.6}{K} = \qty{539.5}{\celsius}$.

\textbf{4.} Penyebut $20.19 + 24.942\times1.6497 = 61.337$; $T = 54496/61.337 =
\qty{888.5}{K} = \qty{615.3}{\celsius}$.

\textbf{5.} $0.311\times1000/67.064 = \qty{4.6}{K}$ untuk inti dan $0.311\times1000/61.337
= \qty{5.1}{K}$ untuk tepi.

\textbf{6.} Termometer berlaku untuk mineral-mineral yang bersama-sama berada dalam kesetimbangan.
Biotit yang terinklusi dalam inti terkurung dari bagian batuan lainnya ketika inti tumbuh;
biotit matriks terakhir kali bertukar dengan tepi. Memasangkan inti dengan matriks berarti mencampur
dua saat yang berbeda.

\noindent\textbf{Bagian II.}

\textbf{7.} Inti $22.80\times539.5 - 1093 = \qty{11208}{bar}$; tepi $22.80\times615.3 -
1093 = \qty{12936}{bar}$.

\textbf{8.} $3\ln(0.050/0.25) = -4.8283$; $RT = 8.314\times812.65 = \qty{6756.4}{J}$; $P =
11208 - 6756.4\times4.8283/6.21 = 11208 - 5253 = \qty{5955}{bar}$.

\textbf{9.} $3\ln(0.060/0.28) = -4.6213$; $RT = 8.314\times888.45 = \qty{7386.6}{J}$; $P
= 12936 - 5497 = \qty{7439}{bar}$.

\textbf{10.} Pada \qty{7439}{bar}: $T = (52630 + 2314)/61.337 = \qty{895.8}{K} =
\qty{622.6}{\celsius}$; pada suhu itu $P = 13102 - 5542 = \qty{7560}{bar}$; lalu $T
= \qty{623.2}{\celsius}$. Tepi tumbuh pada sekitar \qty{623}{\celsius} dan \qty{7.56}{kbar}.

\textbf{11.} Pada \qty{5954}{bar}, $T = \qty{539.2}{\celsius}$ dan $P = \qty{5950}{bar}$:
inti tumbuh pada sekitar \qty{539}{\celsius} dan \qty{5.95}{kbar}. Tebakan pertama,
\qty{6}{kbar}, sudah dekat, dan termometer hanya bergeser \qty{4.6}{K} per kilobar.

\textbf{12.} Inti $5950\times10^5/(2800\times9.81) = \qty{21.7}{km}$; tepi
$7560\times10^5/27468 = \qty{27.5}{km}$.

\textbf{13.} Pada \qty{539}{\celsius} garis itu terletak pada $3.67 + 0.39\times2.10 =
\qty{4.49}{kbar}$, di bawah \qty{5.95}{kbar}: kianit, konsisten dengan kianit yang terinklusi
dalam inti. Pada \qty{623}{\celsius} garis itu terletak pada $5.77 + 0.23\times2.09 = \qty{6.25}{kbar}$,
di bawah \qty{7.56}{kbar}: lagi-lagi kianit.

\noindent\textbf{Bagian III.}

\textbf{14.} $+\qty{84}{\celsius}$ dan $+\qty{1.6}{kbar}$ (sekitar \qty{6}{km} lebih dalam): batuan
itu sedang terkubur dan terpanaskan, pada \omterm{def:b2:pt-paths:paths}{lintasan progradasinya}.

\textbf{15.} Pada \qty{623}{\celsius} atau lebih, di bawah \qty{6.25}{kbar}: setelah tepi tumbuh,
batuan itu kehilangan sekurang-kurangnya \qty{1.3}{kbar}, sekitar \qty{5}{km} beban batuan di atasnya, sambil tetap
sekurang-kurangnya sama panasnya. Batuan itu tersingkap kembali selagi masih panas.

\textbf{16.} Klorit adalah mineral hidrus berderajat rendah: mineral itu merekam \omterm{def:b2:pt-paths:paths}{lintasan retrogradasi}
tahap akhir, pendinginan di kedalaman dangkal, dengan kehadiran air.

\textbf{17.} Penimbunan dan pemanasan sampai tepi, dekompresi tanpa pendinginan, lalu pendinginan:
bila digambar dengan suhu ke kanan dan tekanan ke atas, lintasannya berbelok searah jarum jam.
Puncak tekanan datang lebih dulu, di tepi; puncak termal pada saat yang sama atau lebih kemudian, selama
dekompresi.

\textbf{18.} Ya: model itu juga mencapai puncak tekanannya lebih dulu dan puncak termalnya,
\qty{593}{\celsius} pada \qty{7.8}{kbar}, selama pendakian, dekat dengan kondisi tepi
sekis itu. Penebalan yang diikuti erosi lambat menjelaskan lintasan sekis itu.

\noindent\textbf{Bagian IV.}

\textbf{19.} $(623 - 405)/(412.0 - 398.0) = 218/14.0 = \qty{15.6}{\celsius}$ per juta
tahun.

\textbf{20.} Batuan itu mendingin sekitar satu setengah kali lebih cepat daripada yang diandaikan. Mineral
yang mendingin lebih cepat punya waktu lebih sedikit untuk kehilangan argonnya, sehingga mineral itu menutup pada suhu yang lebih tinggi:
suhu penutupan yang sebenarnya terletak agak di atas \qty{405}{\celsius}, dan penutupannya
agak lebih dalam.

\textbf{21.} $(405 - 10)/30 = \qty{13.2}{km}$.

\textbf{22.} $27.5 - 13.2 = \qty{14.3}{km}$ (dengan kedalaman yang tidak dibulatkan, 27.52 dan 13.17,
\qty{14.35}{km}).

\textbf{23.} $\sqrt{1.5^2 + 1.0^2} = \qty{1.8}{Ma}$: $14.0 \pm \qty{1.8}{Ma}$, jadi lajunya
antara $14.35/15.8 = 0.91$ dan $14.35/12.2 = \qty{1.18}{km/Ma}$.

\textbf{24.} $14.35/14.0 = \qty{1.03}{km}$ per juta tahun, yaitu
\textbf{\qty[text-series-to-math=true]{1.0}{mm.yr^{-1}}}, sekitar tiga puluh kali lebih lambat daripada
tahap pertama \omterm{def:b2:pt-paths:exhumation}{ekshumasi} batuan Alpen terdalam, dan sekitar dua kali laju erosi
model penebalan.
\end{solution}
